\documentclass[11pt]{article}

\usepackage{amsmath,amssymb}
\usepackage{xurl}
\usepackage[colorlinks=true,linkcolor=blue,citecolor=blue,urlcolor=blue]{hyperref}

\input{./command.tex}

\title{Wolf Theorem 6.4: Positive-Map Scope and the Zero-Map Boundary}
\author{}
\date{2026-08-29}

\begin{document}
\maketitle
\gapnote{false-source}{resolved}

\section{Source statement}

Theorem~6.4 of Ref.~\cite{Wolf2012Quantum} concerns a positive linear map
\begin{align}
  T
    &: \MN{D}\longrightarrow\MN{D}
  \label{eq:wolf_thm6_4_positive_map_spectral_characterization_map}
\end{align}
with spectral radius $r$.  It states that irreducibility is equivalent to the
following three properties:
\begin{enumerate}
  \item $r$ is an ordinary nondegenerate eigenvalue of $T$;
  \item there is an $X>0$ satisfying $T(X)=rX$;
  \item there is a $Y>0$ satisfying $T^*(Y)=rY$.
\end{enumerate}
Here $T^*$ is the adjoint for the bilinear trace pairing
\begin{align}
  \tr\!\left(T^*(Y)Z\right)
    &=\tr\!\left(YT(Z)\right).
  \label{eq:wolf_thm6_4_positive_map_spectral_characterization_trace_pairing}
\end{align}
The exact local source is
\path{Notes/WolfNoteTexSource/ch06_spectral_properties.tex}, lines~698--721.
No complete-positivity or Kraus hypothesis occurs in the statement or proof.

The word ``nondegenerate'' is geometric: the ordinary complex eigenspace at
$r$ has dimension one.  The source does not assert algebraic simplicity or a
one-dimensional generalized eigenspace.

\section{Formal statement}

The witness structure \leanid{WolfSpectralProperties} and predicate
\leanid{HasWolfSpectralProperties} in
\doclink{QICLean/Channel/Irreducible/PositiveMapSpectralCharacterization}
record a real number $r>0$
and matrices $X,Y>0$ such that
\begin{align}
  \varrho(T)
    &=r,
  \label{eq:wolf_thm6_4_positive_map_spectral_characterization_radius}\\
  T(X)
    &=rX,
  \label{eq:wolf_thm6_4_positive_map_spectral_characterization_right}\\
  \operatorname{traceAdjointMap}(T)(Y)
    &=rY,
  \label{eq:wolf_thm6_4_positive_map_spectral_characterization_left}\\
  \dim_{\mathbb C}\ker(T-r\Id)
    &=1.
  \label{eq:wolf_thm6_4_positive_map_spectral_characterization_simple}
\end{align}
Thus the left eigenvector is source-facing and does not depend on a choice of
Kraus family.

For a nonzero positive map on a nonzero full matrix algebra, the declaration
\leanid{wolf_theorem_6_4} proves
\begin{align}
  T\text{ is irreducible}
    \quad\Longleftrightarrow\quad
  T\text{ has the properties
    \eqref{eq:wolf_thm6_4_positive_map_spectral_characterization_radius}--
    \eqref{eq:wolf_thm6_4_positive_map_spectral_characterization_simple}}.
  \label{eq:wolf_thm6_4_positive_map_spectral_characterization_iff}
\end{align}
The Lean hypotheses are \leanid{IsPositiveMap T}, $T\neq0$, and
\leanid{NeZero D}.  There is no complete-positivity assumption.

\section{Necessary boundary correction}

The printed statement does not exclude the zero map, although its displayed
condition requires
\begin{align}
  T(X)
    &=rX>0.
  \label{eq:wolf_thm6_4_positive_map_spectral_characterization_printed_strict}
\end{align}
On $\MN{1}$, the zero map is positive and irreducible: the only orthogonal
projections are $0$ and $\Id$.  Its spectral radius is zero, so
\eqref{eq:wolf_thm6_4_positive_map_spectral_characterization_printed_strict}
is impossible.  The printed equivalence is therefore false at this boundary.

The hypothesis $T\neq0$ is the uniform correction.  If $D>1$, irreducibility
already excludes the zero map.  The assumption \leanid{NeZero D} separately
excludes the empty matrix algebra, in which positive-definite density witnesses
are unavailable.

\section{Source-faithful proof}

For the forward implication,
\leanid{hasWolfSpectralProperties_of_irreducible_positive} applies the
positive-map form of Theorem~6.3 to $T$.  It obtains $X>0$, the ordinary
one-dimensional $r$-eigenspace, and the spectral-radius identity.  The
trace-adjoint consequence of Theorem~6.3 supplies $Y>0$ at the same eigenvalue.

For the reverse implication, set $S=\sqrt{Y}$ and define exactly the transform
used at local source lines~712--714,
\begin{align}
  T'
    &=r^{-1}\,\operatorname{similarityMap}(S^{-1},T),
  \qquad
  T'(Z)=r^{-1}S\,T(S^{-1}ZS^{-1})S.
  \label{eq:wolf_thm6_4_positive_map_spectral_characterization_gauge}
\end{align}
Positivity is preserved by congruence and positive rescaling.  Trace
preservation follows directly from
\eqref{eq:wolf_thm6_4_positive_map_spectral_characterization_trace_pairing}:
for $W=S^{-1}ZS^{-1}$,
\begin{align}
  \tr(T'(Z))
    &=r^{-1}\tr(YT(W))
      =r^{-1}\tr(T^*(Y)W)
      =\tr(Z).
  \label{eq:wolf_thm6_4_positive_map_spectral_characterization_tp}
\end{align}
Moreover, $X'=SXS>0$ is fixed by $T'$.

If $T'$ preserved a proper corner $P\MN{D}P$, the positive trace-preserving
Ces\`aro theorem
\leanid{IsPositiveMap.exists_fixed_density_of_preserves_compression} would
produce a stationary density matrix $\tau$ in that corner.  Pulling $\tau$
back through $S^{-1}$ gives an ordinary $r$-eigenvector of $T$.  The
one-dimensional eigenspace therefore makes $\tau$ a scalar multiple of $X'$,
which is impossible because $X'>0$ has full support whereas $\tau$ lies in a
proper corner.  This is precisely the contradiction at local source
lines~715--720.

The earlier finite-Kraus package remains available through
\leanid{HasSpectralProperties}.  The declaration
\leanid{hasSpectralProperties_iff_hasWolfSpectralProperties_of_cp} proves that,
for a nonzero completely positive map, it is equivalent to the source-general
package.  It is therefore a specialization rather than the statement of
Theorem~6.4 itself.

\section{Verdict}

\textbf{Verdict: resolved, with one necessary source correction.}  The former
Kraus/complete-positive restriction and the former uniqueness-only-among-PSD
formulation have been removed from the source-facing theorem.  Ordinary
eigenspace nondegeneracy and the trace-pairing left eigenvector are explicit.
The additional hypothesis $T\neq0$ is necessary only because the printed
equivalence fails for the irreducible zero map on $\MN{1}$.

\bibliographystyle{alpha}
\bibliography{references}

\end{document}
